\chapter{Methodology}
\label{chap:methodology}

This chapter shows a full-width figure, a normal table, and a code listing.

\section{Task and Protocol}
\label{sec:method-protocol}

Quis autem vel eum iure reprehenderit qui in ea voluptate velit esse quam nihil molestiae consequatur, vel illum qui dolorem eum fugiat quo voluptas nulla pariatur.
Scores are reported as macro F1 and recall, with the \ac{aucroc} as context.

\begin{figure}[htbp]
  \centering
  \begin{tikzpicture}[x=1mm, y=1mm, font=\small,
    box/.style={draw=figboxdraw, fill=figbox, rounded corners=1mm, minimum height=9mm, inner sep=2mm},
    arr/.style={-{Latex[length=2mm]}, figboxdraw}]
    \node[box] (a) at (0,0)  {input};
    \node[box] (b) at (30,0) {model};
    \node[box] (c) at (60,0) {output};
    \draw[arr] (a) -- (b);
    \draw[arr] (b) -- (c);
  \end{tikzpicture}
  \caption[A figure in the text block]{A figure at text width. Draw figures with TikZ so they inherit the thesis fonts and colours, which keeps them sharp at any zoom.}
  \label{fig:method-pipeline}
\end{figure}

Ut enim ad minima veniam, quis nostrum exercitationem ullam corporis suscipit laboriosam, nisi ut aliquid ex ea commodi consequatur (\Cref{fig:method-pipeline}).

\section{Models and Training}
\label{sec:method-models}

\begin{table}[htbp]
  \centering
  \caption[A table in the text block]{A normal table. Keep headers as plain text, since an acronym macro would render them as a coloured link.}
  \label{tab:method-config}
  \small
  \begin{tabular}{@{}lr@{}}
    \toprule
    Setting & Value \\
    \midrule
    Hidden width & 256 \\
    Layers & 4 \\
    Learning rate & $10^{-3}$ \\
    Seeds & 42, 43, 44 \\
    \bottomrule
  \end{tabular}
\end{table}

The configuration of \Cref{tab:method-config} stays fixed across the experiments, so the rows of \Cref{chap:results} differ in one factor at a time.

\section{Implementation}
\label{sec:method-implementation}

A listing sets in the monospace font of the template:

\begin{lstlisting}[language=Python, caption={[A code listing]A code listing, numbered and referenced like a figure.}, label=lst:method-example]
def predict(graph, model):
    logits = model(graph.x, graph.edge_index)
    return logits.softmax(dim=1)[:, 1]
\end{lstlisting}

\Cref{lst:method-example} shows the call, and \code{predict} in the text uses the inline command.

\section{Use of Generative AI Tools}
\label{sec:method-ai-use}

% The declaration on the first page points here. The examination office requires the
% product name, the manufacturer, the version, and the purpose of every generative AI
% tool used. Name them in prose and keep the list current.
Generative AI tools were used throughout this thesis and are named here with their purposes.
% Decorative margin illustration (2026-09-08): a robot in the thesis figure palette, no caption
% and no label so it stays out of the list of figures. Units are millimetres, the drawing is
% drawn at 40 mm width and scaled to 60 percent, 24 mm, inside the 41.4 mm margin column.
\begin{marginfigure}
  \centering
  \begin{tikzpicture}[x=0.6mm, y=0.6mm, line join=round,
    part/.style={draw=figboxdraw, fill=figbox, rounded corners=0.7mm},
    joint/.style={draw=figboxdraw, fill=figboxdraw, rounded corners=0.35mm},
    limb/.style={line width=1.6mm, figboxdraw, line cap=round}]
    % legs and feet
    \foreach \x in {-6, 6} {
      \draw[joint] (\x-2.2,3) rectangle (\x+2.2,8);
      \draw[part] (\x-4,0) rectangle (\x+4,3);
    }
    % arms: left hanging, right raised
    \draw[limb] (-13,23) -- (-17,13);
    \draw[limb] (13,23) -- (18,31);
    % body with a screen showing lines of text and a cursor
    \draw[part] (-12,7) rectangle (12,27);
    \draw[draw=figboxdraw, fill=white, rounded corners=0.5mm] (-8,12) rectangle (8,23);
    \foreach \y/\len in {20/10, 17.5/7, 15/5} {
      \draw[figline, line width=0.5mm, line cap=round] (-6,\y) -- (-6+\len,\y);
    }
    \fill[figaccent] (0,14.1) rectangle (1.6,15.9);
    % hands
    \draw[part] (-17,12) circle (2.4);
    \draw[part] (18.4,31.8) circle (2.4);
    % neck, head, ears
    \draw[joint] (-3,27) rectangle (3,30.5);
    \draw[part] (-11,30) rectangle (11,46);
    \draw[joint] (-13,35) rectangle (-11,41);
    \draw[joint] (11,35) rectangle (13,41);
    % eyes with highlights, smile
    \foreach \x in {-4.5, 4.5} {
      \fill[figink] (\x,39) circle (2);
      \fill[white] (\x+0.7,39.7) circle (0.6);
    }
    \draw[figink, line width=0.3mm, line cap=round] (-3.4,35) arc[start angle=-150, end angle=-30, radius=4];
    % antenna with a red tip
    \draw[figboxdraw, line width=0.36mm] (0,46) -- (0,51);
    \fill[figaccent] (0,52.2) circle (1.6);
  \end{tikzpicture}
\end{marginfigure}
The large language models that are the object of study are documented with their versions in \Cref{sec:method-models} and are not meant here.
Name each tool with its product name, its manufacturer, the version you used, and what you used it for, for example writing assistance, language editing, literature search, or code.
Every design decision, every interpretation, and the final wording of every passage are my own, and I take full responsibility for the selection, adoption, and all results of the AI-generated output.
